\documentclass[12pt]{article}

\usepackage{sbc-template}
\usepackage{graphicx,url}
\usepackage[utf8]{inputenc}
\usepackage[brazil]{babel}
\usepackage{amsmath}
\usepackage{amssymb}

\sloppy

\title{Detecção de Comunidades Críticas de Atendimento ao Cliente\\
em Redes Sociais via Busca pelo Clique Máximo}

\author{Matheus Garcia\inst{1}, Luccas Henrique\inst{1}}

\address{Departamento de Ciência da Computação --- Universidade Federal de Roraima (UFRR)\\
  Boa Vista --- RR --- Brasil
  \email{matheusgarciasam@gmail.com, luccas.henrique.vr@gmail.com}
}

\begin{document}

\maketitle

\begin{abstract}
This work presents a combinatorial search engine for detecting critical
customer-support communities in social networks. The interactions among users
are modeled as an undirected graph and the problem of finding densely connected
groups (\emph{crisis cohorts}) is reduced to the \emph{Maximum Clique} problem,
which is NP-hard. Two heuristics were implemented in C11 --- a Greedy
constructive algorithm and a GRASP metaheuristic --- both supported by a bitset
representation that reduces neighborhood intersection to a bitwise AND. We report
a complexity analysis, a theoretical comparison with the exact Bron--Kerbosch
algorithm, and an experimental evaluation over DIMACS benchmarks and a real
Twitter customer-support dataset.
\end{abstract}

\begin{resumo}
Este trabalho apresenta um motor de busca combinatória para a detecção de
comunidades críticas de atendimento ao cliente em redes sociais. As interações
entre usuários são modeladas como um grafo não-direcionado e o problema de
encontrar grupos densamente conectados (\emph{coortes de crise}) é reduzido ao
problema do \emph{Clique Máximo}, NP-difícil. Foram implementadas duas
heurísticas em C11 --- um algoritmo Guloso construtivo e uma metaheurística
GRASP --- ambas apoiadas por uma representação em \emph{bitset} que reduz a
interseção de vizinhanças a um AND bit a bit. Apresentamos a análise de
complexidade, uma comparação teórica com o algoritmo exato de Bron--Kerbosch e
uma avaliação experimental sobre benchmarks DIMACS e um conjunto real de tweets
de atendimento ao cliente.
\end{resumo}

\section{Introdução e Cenário de Aplicação}

Em plataformas de redes sociais, empresas mantêm canais de atendimento ao
cliente que recebem grande volume de interações diárias. Quando ocorre uma falha
sistêmica --- a indisponibilidade de um serviço, uma cobrança indevida em massa
ou um defeito de produto --- surge rapidamente um conjunto de usuários
\emph{altamente interconectados}, todos reagindo ao mesmo problema, respondendo
uns aos outros e mencionando-se mutuamente. Esses agrupamentos, aqui chamados de
\emph{coortes de crise}, são valiosos para a operação de suporte: identificá-los
permite priorizar incidentes sistêmicos, dimensionar equipes e medir o alcance
de uma falha.

O problema de localizar o maior grupo em que \emph{todos} os usuários
interagiram entre si corresponde, em teoria dos grafos, ao problema do
\textbf{Clique Máximo}. Modelando a rede como um grafo não-direcionado
$G = (V, E)$, em que cada vértice é um usuário e cada aresta uma interação
(resposta direta ou menção), uma coorte de crise é um subconjunto
$C \subseteq V$ tal que todo par $(u, v) \in C$ possui aresta direta, e busca-se
maximizar $|C|$. O Clique Máximo é um problema NP-difícil
\cite{garey1979}, o que motiva o emprego de heurísticas eficientes para
instâncias de larga escala.

Este relatório descreve a formulação do problema (Seção~\ref{sec:problema}), os
algoritmos implementados (Seção~\ref{sec:algoritmos}), a análise de complexidade
com comparação teórica ao algoritmo exato de Bron--Kerbosch
(Seção~\ref{sec:complexidade}), os detalhes de implementação
(Seção~\ref{sec:impl}), o estudo comparativo de desempenho
(Seção~\ref{sec:experimentos}) e a aplicação prática sobre dados reais do
Twitter (Seção~\ref{sec:aplicacao}).

\section{Definição Formal do Problema}
\label{sec:problema}

Seja $G = (V, E)$ um grafo não-direcionado simples, com $n = |V|$ vértices e
$m = |E|$ arestas. Um \textbf{clique} é um subconjunto $C \subseteq V$ tal que
$\forall\, u, v \in C,\ u \neq v \Rightarrow (u, v) \in E$, isto é, o subgrafo
induzido por $C$ é completo. O \textbf{Clique Máximo} é o clique de maior
cardinalidade em $G$, e seu tamanho é o \emph{número de clique} $\omega(G)$.

\begin{equation}
\omega(G) = \max\{\, |C| : C \subseteq V \text{ e } C \text{ é clique} \,\}.
\end{equation}

A versão de decisão (``existe clique de tamanho $\geq k$?'') é NP-completa,
sendo um dos 21 problemas clássicos de Karp; a versão de otimização é
NP-difícil. Não se conhece algoritmo polinomial exato, e o problema é, ainda,
difícil de aproximar.

\section{Algoritmos Implementados}
\label{sec:algoritmos}

\subsection{Representação do Grafo}

A escolha da estrutura de dados é determinante para o desempenho da busca por
cliques, cuja operação central é a \emph{interseção de vizinhanças}. O grafo é
representado de forma \textbf{dupla e simultânea}:

\begin{itemize}
  \item \textbf{Bitset (vetor de bits)} --- estrutura principal. Cada vértice
  $i$ possui um vetor de palavras \texttt{uint64\_t} no qual o bit $j$ indica a
  adjacência com o vértice $j$. A interseção de duas vizinhanças --- coração da
  expansão de cliques --- reduz-se a um \emph{AND bit a bit}, percorrendo
  $\lceil n/64 \rceil$ palavras em vez de $n$ posições. A contagem de candidatos
  usa a instrução de hardware \texttt{popcount}.
  \item \textbf{Lista de adjacência} --- estrutura auxiliar para percurso geral
  de vizinhos e navegação esparsa.
\end{itemize}

\subsection{Módulo 1 --- Algoritmo Guloso}

A heurística gulosa constrói um único clique de forma determinística:

\begin{enumerate}
  \item Seleciona como semente o vértice de \emph{maior grau global}.
  \item Mantém um conjunto de candidatos $P$ (em bitset) com os vértices
  adjacentes a todos os membros já no clique.
  \item A cada passo, escolhe de $P$ o candidato com \emph{maior grau dentro de
  $P$} (maior número de vizinhos comuns), maximizando a chance de continuar a
  expansão.
  \item Adiciona o candidato ao clique e atualiza $P \leftarrow P \cap N(v)$ via
  AND de bitsets.
  \item Termina quando $P = \varnothing$.
\end{enumerate}

O algoritmo é rápido, porém míope: por nunca revisar escolhas, fica preso em
ótimos locais.

\subsection{Módulo 2 --- Metaheurística GRASP}

O GRASP (\emph{Greedy Randomized Adaptive Search Procedure}) executa
\texttt{max\_iter} iterações independentes, retendo o melhor clique encontrado.
Cada iteração combina duas fases:

\begin{description}
  \item[Fase construtiva (RCL).] Idêntica ao guloso, exceto que, em vez de
  sempre tomar o melhor candidato, sorteia aleatoriamente entre os
  \texttt{rcl\_size} melhores (a \emph{Restricted Candidate List}). Isso
  diversifica as soluções entre iterações. A aleatoriedade usa um gerador
  congruencial linear (LCG) com semente fixa, garantindo reprodutibilidade.
  \item[Busca local 1-opt.] Tenta substituir um vértice do clique por
  \emph{dois} vértices externos mutuamente adjacentes e adjacentes a todos os
  demais membros. Cada troca bem-sucedida aumenta a cardinalidade em uma
  unidade; o processo repete enquanto houver melhoria.
\end{description}

A combinação de construção aleatorizada com refinamento local permite ao GRASP
escapar dos ótimos locais que limitam o guloso.

\section{Análise de Complexidade}
\label{sec:complexidade}

Seja $n = |V|$, $m = |E|$, $\omega$ o tamanho do clique encontrado e
$W = \lceil n/64 \rceil$ o número de palavras de um bitset. A
Tabela~\ref{tab:complexidade} resume as operações fundamentais.

\begin{table}[ht]
\centering
\caption{Complexidade das operações e algoritmos.}
\label{tab:complexidade}
\begin{tabular}{l l}
\hline
\textbf{Operação / Algoritmo} & \textbf{Complexidade} \\
\hline
AND de bitsets ($P \cap N(v)$)      & $O(n/64)$ \\
\texttt{popcount} de bitset          & $O(n/64)$ \\
Inserção de aresta                   & $O(1)$ \\
Teste de aresta                      & $O(1)$ \\
\hline
Leitor DIMACS                        & $O(n + m)$ \\
Parser Twitter (CSV)                 & $O(L)$ \footnotesize{($L$ = linhas)} \\
\hline
Guloso                               & $O(\omega \cdot n^2 / 64)$ \\
GRASP                                & $O(\texttt{max\_iter} \cdot \omega \cdot n^2 / 64)$ \\
\hline
\end{tabular}
\end{table}

\paragraph{Guloso.} A busca da semente custa $O(n)$. O laço principal executa
até $\omega$ passos; em cada passo varre os candidatos de $P$ (até $n$) e, para
cada um, realiza um AND e um \texttt{popcount} de custo $O(n/64)$. Logo, cada
passo custa $O(n \cdot n/64) = O(n^2/64)$ e o total é
$O(\omega \cdot n^2/64)$. Como $\omega \leq n$, o pior caso é
$O(n^3/64)$, mas na prática $\omega \ll n$.

\paragraph{GRASP.} Cada iteração reproduz o custo da construção gulosa, acrescido
de ordenações $O(n \log n)$ por passo (para montar a RCL) e da busca local, cuja
varredura custa $O(\omega \cdot n^2/64)$ por passada. Multiplicando pelas
\texttt{max\_iter} iterações, obtém-se
$O(\texttt{max\_iter} \cdot \omega \cdot n^2/64)$. O fator $1/64$ proveniente do
bitset é a principal vantagem prática da implementação.

\subsection{Comparação Teórica com o Algoritmo Exato}

Por ser NP-difícil, o Clique Máximo não admite solução exata polinomial
conhecida. Para situar as heurísticas, consideramos dois métodos exatos de
referência.

\paragraph{Força bruta.} Enumerar todos os $2^n$ subconjuntos de $V$ e verificar,
em $O(n^2)$, se cada um é clique resulta em $O(2^n \cdot n^2)$. É inviável já
para $n \approx 40$.

\paragraph{Bron--Kerbosch.} O algoritmo de Bron--Kerbosch
\cite{bron1973} enumera \emph{todos} os cliques maximais por
\emph{backtracking} recursivo sobre três conjuntos: $R$ (clique parcial),
$P$ (candidatos) e $X$ (já explorados):

\begin{verbatim}
BronKerbosch(R, P, X):
    se P = vazio e X = vazio:
        reporta R como clique maximal
    para cada v em P:
        BronKerbosch(R + {v}, P inter N(v), X inter N(v))
        P <- P \ {v}
        X <- X + {v}
\end{verbatim}

Com a escolha de um \emph{pivô} para reduzir a ramificação, sua complexidade de
pior caso é $O(3^{n/3})$, que é ótima no sentido de Moon--Moser: um grafo de
$n$ vértices pode ter até $3^{n/3}$ cliques maximais. Com ordenação por
\emph{degeneracy} $d$, obtém-se $O(d \cdot n \cdot 3^{d/3})$, eficiente para
grafos esparsos. O Clique Máximo seria então o maior clique maximal reportado.

A Tabela~\ref{tab:exato} contrasta as abordagens. Enquanto Bron--Kerbosch
\emph{garante} a solução ótima, seu custo exponencial o restringe a instâncias
pequenas; as heurísticas trocam essa garantia por escalabilidade polinomial.

\begin{table}[ht]
\centering
\caption{Heurísticas vs.\ algoritmo exato.}
\label{tab:exato}
\begin{tabular}{l c c}
\hline
\textbf{Método} & \textbf{Complexidade} & \textbf{Garante ótimo?} \\
\hline
Força bruta        & $O(2^n \cdot n^2)$        & Sim \\
Bron--Kerbosch     & $O(3^{n/3})$             & Sim \\
Guloso             & $O(\omega \cdot n^2/64)$ & Não \\
GRASP              & $O(\texttt{max\_iter}\cdot\omega\cdot n^2/64)$ & Não \\
\hline
\end{tabular}
\end{table}

\section{Detalhes de Implementação}
\label{sec:impl}

O projeto foi desenvolvido inteiramente em \textbf{C (padrão C11)}, sem
bibliotecas externas de manipulação de grafos, e organizado em módulos coesos:

\begin{itemize}
  \item \texttt{graph.[ch]} --- estrutura do grafo (bitset + lista de
  adjacência), criação, inserção/teste de arestas e estatísticas.
  \item \texttt{clique.[ch]} --- algoritmos Guloso e GRASP, e a estrutura
  \texttt{CliqueResult} (vértices, tamanho, iterações, tempo).
  \item \texttt{dimacs.[ch]} --- leitor de benchmarks no formato DIMACS
  \texttt{.clq} (linhas \texttt{p edge} e \texttt{e u v}).
  \item \texttt{twitter\_parser.[ch]} --- parser de CSV do dataset Twitter, com
  tabela hash \texttt{string}$\rightarrow$\texttt{int} (\emph{djb2}) para mapear
  identificadores e resolver as arestas de resposta em passagem única.
  \item \texttt{dot\_export.[ch]} --- exportação para o formato GraphViz
  \texttt{.dot}, destacando o clique em vermelho.
  \item \texttt{main.c} --- ponto de entrada com os modos \texttt{test},
  \texttt{greedy}, \texttt{grasp}, \texttt{twitter} e \texttt{bench}.
\end{itemize}

A compilação usa \texttt{make} com as flags
\texttt{-Wall -Wextra -O2 -std=c11}; o projeto compila sem nenhum aviso. A
gestão de memória é explícita, com liberação de todas as estruturas
(\texttt{graph\_free}, \texttt{clique\_result\_free}, \texttt{hashmap\_free}).

\section{Estudo Comparativo de Desempenho}
\label{sec:experimentos}

\subsection{Metodologia}

Os experimentos usaram dois benchmarks padrão da competição DIMACS, amplamente
empregados na literatura de Clique Máximo. O GRASP foi configurado com
\texttt{max\_iter}~$=100$, \texttt{rcl\_size}~$=3$ e semente~$=42$. A qualidade
é medida como a razão entre o clique obtido e o ótimo conhecido. Os tempos foram
obtidos via \texttt{clock()} (CPU time), com binário compilado em \texttt{-O2}.
A Tabela~\ref{tab:grafos} caracteriza as instâncias.

\begin{table}[ht]
\centering
\caption{Características das instâncias DIMACS.}
\label{tab:grafos}
\begin{tabular}{l c c c c c}
\hline
\textbf{Instância} & \textbf{Vért.} & \textbf{Arestas} & \textbf{Densidade}
 & \textbf{Grau máx.} & \textbf{$\omega$} \\
\hline
\texttt{brock200\_2}  & 200 & 9876  & 0{,}4963 & 114 & 12 \\
\texttt{p\_hat300-1}  & 300 & 10933 & 0{,}2438 & 132 & 8  \\
\hline
\end{tabular}
\end{table}

\subsection{Resultados}

A Tabela~\ref{tab:resultados} apresenta os resultados. O número de arestas lido
coincidiu exatamente com o esperado em ambos os arquivos, validando o leitor
DIMACS.

\begin{table}[ht]
\centering
\caption{Comparativo Guloso vs.\ GRASP nos benchmarks DIMACS.}
\label{tab:resultados}
\begin{tabular}{l l c c c c}
\hline
\textbf{Instância} & \textbf{Algoritmo} & \textbf{Clique} & \textbf{Ótimo}
 & \textbf{Qualidade} & \textbf{Tempo (ms)} \\
\hline
\texttt{brock200\_2} & Guloso & 9  & 12 & 75{,}0\,\%  & 0{,}39  \\
\texttt{brock200\_2} & GRASP  & 10 & 12 & 83{,}3\,\%  & 49{,}85 \\
\hline
\texttt{p\_hat300-1} & Guloso & 7  & 8  & 87{,}5\,\%  & 0{,}34  \\
\texttt{p\_hat300-1} & GRASP  & 8  & 8  & 100{,}0\,\% & 42{,}04 \\
\hline
\end{tabular}
\end{table}

\subsection{Discussão}

Os resultados confirmam o \emph{trade-off} esperado entre qualidade e tempo:

\begin{itemize}
  \item O \textbf{GRASP supera o Guloso} em ambas as instâncias, atingindo o
  \textbf{ótimo} ($\omega = 8$) em \texttt{p\_hat300-1} e ficando a apenas dois
  vértices do ótimo em \texttt{brock200\_2}, instância notoriamente difícil.
  \item O custo é maior: o GRASP leva $\sim$40--50\,ms contra $\sim$0{,}3\,ms do
  Guloso (cerca de 100$\times$), coerente com o fator \texttt{max\_iter}~$=100$.
  \item O Guloso é apropriado para uma primeira estimativa em tempo quase
  instantâneo; o GRASP é preferível quando a qualidade da solução é prioritária
  e o orçamento de tempo permite.
\end{itemize}

A diferença em \texttt{brock200\_2} ilustra a dificuldade da família
\emph{brock}, projetada com cliques ``escondidos'' que resistem a heurísticas
locais; mais iterações ou uma busca local mais agressiva (2-opt) poderiam
reduzir essa lacuna.

\section{Aplicação Prática: Dataset Twitter}
\label{sec:aplicacao}

Para demonstrar a aplicação no cenário-alvo, o módulo Twitter processou uma
amostra real do dataset \emph{Customer Support on Twitter} (Kaggle). O parser
constrói o grafo a partir do campo \texttt{in\_response\_to\_tweet\_id},
ligando autores que respondem uns aos outros. Na amostra de 99 tweets, obteve-se
um grafo com 44 autores, 28 arestas e densidade 0{,}0296; a busca gulosa
identificou uma coorte (clique) de tamanho 2, compatível com a baixa densidade
de uma amostra pequena. Sobre o dataset completo ($\sim$3 milhões de tweets), a
mesma abordagem escala graças à passagem única do parser ($O(L)$) e à
representação em bitset. O resultado é exportado em \texttt{.dot} e renderizado
com GraphViz, destacando a coorte de crise em vermelho.

\section{Conclusão}

Este trabalho modelou a detecção de comunidades críticas de atendimento ao
cliente como o problema do Clique Máximo e implementou, em C11, duas heurísticas
--- Guloso e GRASP --- apoiadas por uma representação em bitset que torna a
interseção de vizinhanças um AND bit a bit. A análise de complexidade situou as
heurísticas como soluções polinomiais frente ao custo exponencial dos métodos
exatos (força bruta e Bron--Kerbosch). Experimentalmente, o GRASP alcançou o
ótimo em uma das instâncias e superou consistentemente o Guloso, ao custo de
maior tempo de execução. A aplicação ao dataset real do Twitter demonstrou a
viabilidade prática da abordagem. Como trabalho futuro, destacam-se a
implementação de uma busca local 2-opt, o ajuste adaptativo de
\texttt{max\_iter}/\texttt{rcl\_size} e a comparação empírica direta com o
Bron--Kerbosch em instâncias pequenas.

\section*{Repositório}

O código-fonte completo, os benchmarks e as instruções de compilação estão
disponíveis em:
\begin{center}
\url{https://github.com/Luccas-H/FP_DCC606_Tema_6_RR_2026}
\end{center}

\begin{thebibliography}{9}

\bibitem{garey1979}
GAREY, M. R.; JOHNSON, D. S.
\emph{Computers and Intractability: A Guide to the Theory of NP-Completeness}.
W. H. Freeman, 1979.

\bibitem{bron1973}
BRON, C.; KERBOSCH, J.
Algorithm 457: Finding All Cliques of an Undirected Graph.
\emph{Communications of the ACM}, v.~16, n.~9, p.~575--577, 1973.

\bibitem{santiago2018}
BEHAR, R.; COHEN, S.
Finding All Maximal Connected $s$-Cliques in Social Networks.
In: \emph{Proceedings of the 21st International Conference on Extending Database
Technology (EDBT)}, p.~61--72, 2018. DOI: 10.5441/002/edbt.2018.07.

\bibitem{feo1995}
FEO, T. A.; RESENDE, M. G. C.
Greedy Randomized Adaptive Search Procedures.
\emph{Journal of Global Optimization}, v.~6, p.~109--133, 1995.

\bibitem{dimacs}
DIMACS Implementation Challenge --- Maximum Clique.
Disponível em: \url{http://archive.dimacs.rutgers.edu/Challenges/}.

\end{thebibliography}

\end{document}
